% =====================================================================
% 10c-appendix-tool.tex -- T-UCA full table, fields, reproduction
% SOURCE OF TRUTH: docs/24 (tool-gateway UCA contract, v1.1: T-UCA-01..16,
% policy_version "t-uca-1.1"; audit field for tool-side ids is `ucas`).
% PENDING NUMBERS (fill only from a5 real outputs): N/M/K/A/X in 06e.
% =====================================================================

\subsection{Tool-Call Unsafe Control Actions ({T-UCA})}\label{app:tuca}

This appendix mirrors the tool-side matrix that the implementation encodes. The
matrix is defined once, in the project's tool-gateway contract
(\texttt{docs/24}), and is implemented twice---as \texttt{engine/\allowbreak{}tool\_\allowbreak{}gateway.\allowbreak{}py}
(pure functions, zero dependencies) and as
\texttt{plugin/\allowbreak{}dsh-\allowbreak{}lingshu/\allowbreak{}tool-\allowbreak{}gateway.\allowbreak{}js} (same semantics)---with both rule
tables machine-synchronised from the contract's rule array, so that each
identifier, verdict and hint text agrees row by row. The namespace is disjoint
from the injection-side UCA-01--08; the two are related only by ordering:
injection pre-flight runs before the model call, the tool gate runs after a call
is emitted and before execution. The contract's policy version is
\texttt{t-uca-1.1} (the \texttt{v1.0.0-paper-baseline} snapshot, commit
\texttt{d551079}, 2026-10-27), and the field carrying tool-side identifiers in the audit
event is named \texttt{ucas} (same spelling as the injection-side field, but a
different event type: \texttt{tool\_gate.ucas} may contain only
\texttt{T-UCA-xx}, \texttt{audit.ucas} only \texttt{UCA-0x}).

\begin{table*}[t]
\centering
\caption{Tool-side unsafe control actions (\texttt{docs/24} v1.1, policy
\texttt{t-uca-1.1}, paper baseline). Sixteen rules: thirteen block, three degrade (a degrade
falls back to block when no approval channel exists). No rule returns
\emph{allow}; allow is the default policy for read-only calls that match no rule.
The \emph{group} column follows the contract's grouping and is not a runtime
field.}
\label{tab:tuca}
\scriptsize
\begin{tabular}{llp{3.6cm}lp{2.2cm}l}
\toprule
ID & Group & Trigger (contract summary) & State inputs & Verdict & Test case \\
\midrule
T-UCA-01 & state        & mutating call while variety fallback or unresolved conflict is active & variety, conflict & block & \texttt{test\_tuca\_01} \\
T-UCA-02 & data         & credential or secret material in arguments or target & capabilities & block & \texttt{test\_tuca\_02} \\
T-UCA-03 & mutation     & destructive delete/overwrite without exact matching confirmation & --- & block & \texttt{test\_tuca\_03} \\
T-UCA-04 & mutation     & write to frozen/protected path (incl.\ \texttt{engine/}, \texttt{plugin/} sources) & frozen\_ok & block & \texttt{test\_tuca\_04} \\
T-UCA-05 & network/exec & network egress outside the (default empty) allowlist & capabilities & block & \texttt{test\_tuca\_05} \\
T-UCA-06 & network/exec & arbitrary code execution (capability or execution-shaped args) & capabilities & block & \texttt{test\_tuca\_06} \\
T-UCA-07 & shape        & oversized, deeply nested or multi-encoded payload & --- & block & \texttt{test\_tuca\_07} \\
T-UCA-08 & shape        & unregistered tool or incomplete registry entry & tool\_registry & block & \texttt{test\_tuca\_08} \\
T-UCA-09 & shape        & schema violation, delimiter/control-character injection, homoglyph tool name & tool\_registry & block & \texttt{test\_tuca\_09} \\
T-UCA-10 & state        & mutating call while context phase is critical & phase & degrade$\rightarrow$block & \texttt{test\_tuca\_10} \\
T-UCA-11 & audit        & incomplete audit context (missing session/turn/call id) & audit\_context & block & \texttt{test\_tuca\_11} \\
T-UCA-12 & shape        & capability neither declared nor granted & capabilities & block & \texttt{test\_tuca\_12} \\
T-UCA-13 & state        & saturated affective baseline with irreversible write, unconfirmed & saturated & degrade$\rightarrow$block & \texttt{test\_tuca\_13} \\
T-UCA-14 & mutation     & frozen-region check failed and call is mutating & frozen\_ok & block & \texttt{test\_tuca\_14} \\
T-UCA-15 & mutation     & path escapes the workspace (\texttt{..}, absolute, drive, UNC, null byte) & --- & block & \texttt{test\_tuca\_15} \\
T-UCA-16 & state        & destructive call while context phase is shift & phase & degrade$\rightarrow$block & \texttt{test\_tuca\_16} \\
\bottomrule
\end{tabular}
\end{table*}

\paragraph{Default policy and evaluation order.} Unknown, unregistered or
undecidable calls are blocked (\mbox{fail-closed}); read-only calls that match no rule
are allowed and logged; a mutating call whose phase is unknown is blocked; a
parse failure is blocked with an explicit \texttt{parse\_error} reason rather
than inferred intent. When several rules match, the verdict takes the highest
severity (block over degrade over allow) and \emph{all} matched identifiers are
recorded in \texttt{ucas} without truncation. Evaluation order puts static checks
first, then the state-dependent rules, and the audit rule last:
\texttt{T-\allowbreak{}UCA-\allowbreak{}01, 02, 03, 04, 05, 06, 07, 08, 09, 12, 14, 15, 10, 13, 16, 11}.
A \texttt{degrade} means the original call is not executed: the caller returns a
strong hint, sets \texttt{requires\_\allowbreak{}confirmation}, and requires an explicit human
confirmation through the host approval seam; a degrade with no approval channel
is treated as a block (\texttt{degraded\_to}). The gate never executes a payload
and never rewrites arguments; it classifies, consults state, and returns a
verdict.

\paragraph{State inputs and tool classes.} The gate consumes only state that
already exists in the system: \texttt{variety\_fallback} and
\texttt{unresolvedConflict} (from the conflict and variety machinery),
\texttt{phase} (from the preemptive monitor), \texttt{saturated} (from the
integral baseline), \texttt{frozen\_ok} (from the frozen-region check),
\texttt{capabilities} (declared versus granted), \texttt{audit\_context}, and the
tool registry. Tool classes are closed and static---\texttt{read}, \texttt{write},
\texttt{destructive}, \texttt{exec}, \texttt{network}---with \texttt{write},
\texttt{destructive}, \texttt{exec} and \texttt{network} treated as mutating and
\texttt{read} as read-only; a payload cannot declare its own class. Finer
semantics (for example which arguments name a path) live in the registry's
capability declarations rather than in the class enum.

\paragraph{Audit fields.} Each verdict produces one \texttt{tool\_gate} event:
\texttt{schema}, \texttt{policy\_version}, \texttt{call\_id},
\texttt{session\_id}, \texttt{turn\_id}, \texttt{tool}, \texttt{tool\_class},
\texttt{verdict}, \texttt{allowed}, \texttt{ucas}, \texttt{degraded\_to},
\texttt{requires\_\allowbreak{}confirmation}, \texttt{reason}, a nested
\texttt{audit} block (\texttt{args\_digest}, \texttt{args\_keys},
\texttt{args\_bytes}, \texttt{args\_preview}, \texttt{phase},
\texttt{variety\_fallback}, \texttt{frozen\_ok}, \texttt{saturated}), and a
\texttt{snapshot} of the target (path, modification time, hash) used to check
that the object judged is the object executed. Argument digests and key lists
are recorded instead of raw secret-bearing arguments. In the host plugin the
same event drives a clarification flag that surfaces as a strong hint on the next
turn.

\paragraph{Reproduction.} \texttt{python tests/\allowbreak{}test\_\allowbreak{}tool\_\allowbreak{}gateway.\allowbreak{}py} exercises
dangerous, benign, state-gated and cross-language cases (as run for this draft:
79 passing assertions, 20/20 dangerous blocked, 0/8 benign blocked, 20/20
cross-language agreement);
\texttt{node plugin/\allowbreak{}dsh-\allowbreak{}lingshu/\allowbreak{}tool-\allowbreak{}gateway-\allowbreak{}sim.\allowbreak{}mjs -\allowbreak{}-\allowbreak{}assert} checks the rules
individually in JavaScript (65 checks);
\texttt{python tools/\allowbreak{}run\_\allowbreak{}actuator\_\allowbreak{}demo.\allowbreak{}py} drives the reference loop and emits
the JSON cited in Section~\ref{sec:actuator}. The frozen blind held-out set is
\texttt{paper/\allowbreak{}ablation/\allowbreak{}tool/\allowbreak{}heldout-\allowbreak{}tool-\allowbreak{}calls-\allowbreak{}v1.\allowbreak{}jsonl} (ten calls, two per
family; SHA-256
\texttt{\seqsplit{f68fc1b3f19fc37e43b9856096d29b9a55de7dee9bf664a9b9a9592e070576d4}}); it
is evaluated by a party other than its author, under two regimes (arguments
only, and arguments with context), because the conflict family depends on
context.

\paragraph{Boundaries recorded with the matrix.} The host exposes no
pre-dispatch hook, so plugin-side blocking is limited to tools routed through
\texttt{guardTool}; engine-side pre-blocking is a reference loop, not a product
enforcement point. Tools executed outside the gateway are not gated. The gate is
a static checker, so parsing is its own attack surface, and its false-positive
rate is uncalibrated in the absence of real traffic. The contract also records
known false-positive surfaces (confirmation fatigue for writes during the
critical phase, legitimate maintenance of protected paths, documents that merely
discuss credential formats, and large legitimate imports). These are repeated
here because a matrix in an appendix is exactly where a reader is most likely to
mistake a rule table for an enforcement guarantee.
